%! Departures from Linearity — Unit 2, Lesson 2.5 — Lesson Plan
% GRADUAL-RELEASE plan (warm-up / guided notes and practice / debrief / close and start the
% homework) at the COMPACT budget: the whole lesson is finished in one 55-minute period.
% THERE IS NO GROUP ACTIVITY IN THIS COURSE and NO INDEPENDENT PRACTICE SET in the notes — the
% homework is the individual practice, started in class. Teacher-facing. Every box is `skillbox`
% (breakable) with a \boxguard ahead of it; no `fixedskillbox`, no tier boxes, no exit ticket.
% The last lesson of Unit 2: the homework preview points at the unit test and Unit 3.
\documentclass[10pt]{article}
\usepackage{apstats-article}
\usepackage{apstats-boxes}

\newcommand{\UnitNumberName}{Unit 2: Exploring Two-Variable Data \quad}
\newcommand{\LessonNumberName}{Lesson 2.5: Departures from Linearity}

\begin{document}
\begin{center}
    {\Large\bfseries \CourseName \\
    {\small\bfseries \UnitNumberName \LessonNumberName}}
\end{center}

\begin{tcolorbox}[colback=sky, colframe=navy, boxrule=0.9pt, arc=2mm,
    left=2mm, right=2mm, top=1.5mm, bottom=1.5mm]
    \textbf{Primary Objective:} Students will classify an unusual point in a regression as an
    outlier, a high-leverage point, or both, decide from regression output whether it is
    influential, use a residual plot --- not $r$ --- to judge whether a linear model is
    appropriate, and make a prediction from a log-transformed least-squares line.
    \\[2pt]
    \textbf{CED alignment:} Topic 2.9 (with Topic 2.7 carried forward) \quad$\cdot$\quad Big
    Idea: DAT (Data Analysis) \quad$\cdot$\quad Skills 2.A, 2.C \quad$\cdot$\quad LO DAT-1.I,
    DAT-1.J (DAT-1.F) \quad$\cdot$\quad EK DAT-1.I.1--1.I.3, DAT-1.J.1, 1.J.2, DAT-1.F.1, 1.F.2
    \\[2pt]
    \textbf{Lesson model:} \emph{Gradual release --- I do, we do, you do --- all of it inside
    the guided notes.} There is no group activity. The notes name the vocabulary as they teach
    it and deliver the instruction row by row in one \emph{Main Ideas / Notes} table --- each
    idea a definition to complete and a grid of short problems, the first worked by you and the
    rest by the class; the last row, Guided Practice, is worked together with students holding
    the pen; \textbf{the homework is the individual practice}, and students start it in class
    while you circulate. The whole-class debrief is \emph{spoken}.
    \\[2pt]
    \textbf{Compact budget:} the packet is six pages (cover, warm-up, two pages of notes, one
    of AP Practice, one of homework) so the whole lesson finishes in the period. Every
    scatterplot and residual plot is \textbf{given}; students read and annotate, never sketch.
\end{tcolorbox}

\renewcommand{\arraystretch}{1.4}
\boxguard
\begin{skillbox}[Learning Targets \& Key Understandings]{goldbox}
\begin{tabularx}{\linewidth}{@{}X|X@{}}
\textbf{Learning targets (student language)} & \textbf{Key understandings (the ``why'')} \\[2pt]
\hline
\vspace{-6pt}
\begin{itemize}[leftmargin=1.1em, nosep, after=\vspace{2pt}]
  \item I can tell a point that is far from the trend from a point that is far out in $x$.
  \item I can decide from ``with'' and ``without'' output whether a point is influential.
  \item I can use a residual plot to decide whether a line is the right model, whatever $r$
        says.
  \item I can predict from a line fit to $\log y$ and explain why that model fits better.
\end{itemize}
&
\vspace{-6pt}
\begin{itemize}[leftmargin=1.1em, nosep, after=\vspace{2pt}]
  \item An outlier has a large residual; a high-leverage point has an extreme $x$-value; an
        influential point changes the slope, intercept, or $r$ substantially when removed.
        Outliers and high-leverage points are \emph{often}, not always, influential
        (EK DAT-1.I.1--1.I.3).
  \item Random scatter in a residual plot is evidence of linear form; a pattern says the model
        is inappropriate (EK DAT-1.F.1, 1.F.2).
  \item \textbf{Target misconception:} ``$r$ is close to $\pm1$, so the linear model is
        appropriate.'' $r$ measures the \emph{strength} of a linear association; it does not
        certify the \emph{form}. The antibiotic data give $r = -0.94$ and a U-shaped residual
        plot --- a strong curve, and the wrong model.
  \item \textbf{Second misconception:} ``a point far off the trend always drags the line.''
        Shop A sits far above the trend and barely moves the slope; shop C, far out in $x$,
        cuts it by three quarters.
  \item A transformation such as $\log y$ can make curved data linear; more random residuals
        and $r$ closer to $\pm1$ are the evidence it worked (EK DAT-1.J.1, 1.J.2).
\end{itemize}
\end{tabularx}
\end{skillbox}

\boxguard
\begin{skillbox}[{Vocabulary, Concepts, \& Theorems --- taught directly in the guided notes}]{greenbox}
\begin{minipage}[t]{0.48\linewidth}
    \begin{itemize}
        \item \textbf{Outlier (in regression)} --- does not follow the trend; large residual
              (EK DAT-1.I.1)
        \item \textbf{High-leverage point} --- $x$-value much larger or smaller than the rest
              (EK DAT-1.I.2)
        \item \textbf{Influential point} --- removing it changes the slope, intercept, and/or
              $r$ substantially (EK DAT-1.I.3)
    \end{itemize}
\end{minipage}
\hfill
\begin{minipage}[t]{0.48\linewidth}
    \begin{itemize}
        \item \textbf{Transformation} --- e.g.\ $\log y$; named in the Residual Plots row, not
              in the vocabulary box (EK DAT-1.J.1)
        \item \textbf{Predicting from a log model} --- $\hat{y} = 10^{\,a + bx}$ from
              $\widehat{\log y} = a + bx$
        \item \textbf{Carried forward} --- residual $= y - \hat{y}$ and residual plots
              (Lesson 2.3); $r$, $r^2$, and slope interpretation (Lessons 2.2, 2.4)
    \end{itemize}
\end{minipage}
\end{skillbox}

% A tabularx never splits, so guard generously ahead of it.
\boxguard[30]
\begin{skillbox}[Lesson at a Glance --- \MeetingLength]{sky}
\renewcommand{\arraystretch}{1.25}
\begin{tabularx}{\linewidth}{@{}l c X X@{}}
\textbf{Phase} & \textbf{Min} & \textbf{Students} & \textbf{Teacher} \\
\hline
Warm-Up & 5 & A taxi-fare residual, two residual plots, and P or Q & Debrief item~3 aloud;
leave ``Q --- it's far out in $x$'' on the board \\
Guided Notes \& Practice & 34 & Fill the vocabulary box and the notes table: problems 1--9 row
by row, then \emph{The Lemonade Stand} (10--13) together, pen in their hand & \textbf{I do}
each row's definition lines and first problem (1, 5, 8), \textbf{we do} the rest (24)
$\to$ \textbf{we do} Guided Practice (10) \\
Debrief & 8 & Pens down, papers up --- answer aloud, nothing to fill in & Put $r = -0.94$
(wrong model) beside $r = 0.80$ (right form); open on the almost-right problem~7 \\
Close \& Assign & 8 & \textbf{Start the homework in class}, alone and silent & Circulate on
Part~B item~1 for your formative read; preview the Unit 2 test and Unit 3 \\
\end{tabularx}
\end{skillbox}

\boxguard[20]
\begin{skillbox}[Warm-Up --- Activate Prior Knowledge (5 min)]{sky}
\begin{minipage}[t]{0.48\linewidth}
    \textbf{The three items and what each seeds}
    \begin{itemize}
        \item \textbf{1.} A taxi fare: predicted \$10.50, residual $+\$1.50$, too low.
              \textbf{Spirals} to Lessons 2.3--2.4; the notes lean on residuals all period.
        \item \textbf{2.} Two residual plots, one random and one a U. \textbf{Seeds} the
              Residual Plots row: students say ``B curves, so no line'' here, with no $r$ in
              sight --- then problem~7 puts $r = -0.94$ next to the same U.
        \item \textbf{3.} Point P (mid-$x$, above the line) or Q (far right, below). \textbf{Seeds}
              leverage and influence in the Unusual Points row.
    \end{itemize}
\end{minipage}
\hfill
\begin{minipage}[t]{0.48\linewidth}
    \textbf{Running it}
    \begin{itemize}
        \item \textbf{Debrief item~3 aloud} and write the room's reason on the board in their
              words (``Q is way out to the side''). Row~1 names that idea \emph{high leverage}
              ten minutes later.
        \item Expect votes for P on item~3 --- it looks more ``wrong.'' Do not settle it: tell
              them problem~3 in the notes decides it with numbers.
        \item Item~1 is mechanical; take 60 seconds on it and move on.
    \end{itemize}
\end{minipage}
\end{skillbox}

\boxguard
\begin{skillbox}[Guided Notes \& Practice --- I do, then we do (34 min)]{sky}
\textbf{This is the whole lesson.} There is no group activity, and there is no independent
practice set on this page: \textbf{the homework is the individual practice}, started in class.
The notes are one \emph{Main Ideas / Notes} table --- three rows, 13 short problems.
\textbf{In every row you fill the definition lines and work the first problem of each grid; the
class works the rest, pen in their hand.} The vocabulary box is filled \emph{as you go}.
\begin{multicols}{2}
\textbf{Unusual Points (10 min).} The three definitions, then the coffee-shop plot and its
four fitted lines. Work \textbf{problem~1} yourself (A: typical seats, far above the trend ---
an outlier, not high-leverage). Problems 2--4 are theirs. \textbf{Problem~3 is the trap}: A is
the most visibly ``wrong'' point, and adding it moves the slope only from 4.14 to 4.10 --- it
lowers $r$, nothing more. Problem~4 lands C as the influential one: slope 4.14 $\to$ 1.01,
$r$ 0.97 $\to$ 0.48. Point back at P and Q on the board.

\textbf{Residual Plots (14 min) --- the crux.} Put the scatterplot up first and ask how good
the line looks; let them admire $r = -0.94$. Then the residual plot. Work \textbf{problem~5}
(a U: positive at both ends, negative in the middle); problem~6 is theirs.
\textbf{Problem~7 is the crux}: take a hand count on ``agree'' \emph{before} anyone explains.
Then write the callout --- $r$ measures strength, the residual plot checks \emph{form} ---
and name the transformation. Work \textbf{problem~8} with them (two pieces of evidence);
problem~9 is the log prediction, $10^{0.71} \approx 5.1$ mg/L. Keep the transformation light:
read it from output, predict once, move on.

\columnbreak
\textbf{We do (about 10 min): The Lemonade Stand.} A fresh context, problems 10--13, worked
entirely together. Students hold the pen; you hold the questions: \emph{``Show me a curve in
that plot.''} \quad \emph{``Where is 82\textdegree F on the $x$-axis --- the middle or the
edge?''} \textbf{Problem~11 is the crux transferred in reverse}: $r = 0.80$ with random
residuals, so the line is \emph{appropriate} and merely moderate. It is the one to protect
when time is short --- a student who can answer both 7 and 11 has separated strength from form.

\textbf{The pen must actually change hands.} Put problem~10 on them cold, take a hand count on
problem~11 before anyone speaks, and make the room defend both answers. Circulate and demand
the reason:
\begin{itemize}[leftmargin=1.1em, nosep]
  \item \textbf{Problem~11:} ``$r$ went down. Did the \emph{shape} change?''
  \item \textbf{Problem~12:} ``Is 82\textdegree F unusual? Then is it high-leverage?''
  \item \textbf{Problem~13:} ``Which number moved --- the slope or $r$?''
\end{itemize}
While circulating, pick the papers to put up in the debrief --- one \emph{almost}-right
problem~7 (``disagree, because $r$ isn't 1'') and one that names the curve.
\textbf{Your formative read comes twice}: here, and again in the last eight minutes as students
start the homework.
\end{multicols}
\end{skillbox}

\boxguard
\begin{skillbox}[Debrief --- whole class (8 min)]{sky}
Pens down, papers up. \textbf{This debrief is spoken} --- nothing on the page to fill in.
\begin{multicols}{2}
\begin{enumerate}[leftmargin=1.3em, nosep, itemsep=3pt]
  \item \textbf{Open on the almost-right problem~7}: ``disagree, because $r$ is not exactly
        $-1$.'' Right verdict, wrong reason --- by that logic $r = -0.99$ would rescue the line.
        Correct it gently and in public: the reason is the U.
  \item Put the two pairs side by side on the board: antibiotic, $r = -0.94$, \emph{wrong}
        model; lemonade, $r = 0.80$, \emph{right} form. That board is the lesson.
  \item Take problem~3 back: why did A barely move the line while C dragged it? Leverage.
\end{enumerate}
\columnbreak
\textbf{The sentence to leave on every desk:} \emph{$r$ tells you how strong; the residual plot
tells you whether a line is the right shape at all.} Make a student say it, not you.

\textbf{If the debrief runs short}, cut item~3 --- item~1 is the only one that cannot be cut.
\textbf{Do not let the debrief eat the last eight minutes} --- the homework start is your
second formative read.
\end{multicols}
\end{skillbox}

\boxguard
\begin{skillbox}[AP Practice --- assigned, not scored]{goldbox}
Two multiple-choice items on a bacteria culture --- \textbf{MC~1 is the crux in AP clothing}
($r = 0.96$ and a curved residual plot; distractor (A) is the target misconception word for
word) and MC~2 is a log-model prediction ($10^{1.5} \approx 31.6$; distractor (A) forgets to
undo the log) --- then one free-response set on 13 used cars: \textbf{(a)} is the
210{,}000-mile car high-leverage, \textbf{(b)} is it influential (slope $-0.103 \to -0.055$),
and \textbf{(c)}, \textbf{the spiral}, interpreting the slope in context (Lesson~2.4).

\textbf{Not scored --- the cover's score column reads \textbf{NA} for this page.} It is also the
best review for the unit test; go over it in the review class. Hold the AP standard out loud:
\emph{in context or it does not count} --- a part (c) that says ``the price goes down 0.103''
with no units and no ``predicted'' has not earned the point.
\end{skillbox}

\boxguard
\begin{skillbox}[Homework --- scored, due the first class after two study halls]{goldbox}
Five items in a third context. \textbf{Part~A} (apartment rents): \textbf{A1} is P high-leverage
and influential, from with/without output; \textbf{A2} an AP-style multiple choice (a far-out
point \emph{on} the line barely moves the slope) with a required justification. \textbf{Part~B}
(phone resale values): \textbf{B1} the crux restated --- $r = -0.94$ beside a U-shaped
residual plot; \textbf{B2} a log-model prediction ($10^{2.24} \approx \$174$); \textbf{B3} two
pieces of evidence that the log model fits better.

\textbf{Start it in the last eight minutes of the period, in class and alone.}
\textbf{Scoring:} 10 points --- 2 per item. \textbf{B1 is the one to read closely}: a student
who agrees with the classmate has kept the target misconception, and it will cost them on the
unit test's free response. The spiralbox previews the Unit~2 test and Unit~3 instead of a next
lesson.

\textbf{Paper homework is the default for this course} --- assign this page unless you have
decided otherwise. \textbf{DeltaMath override (occasional):} on the days you assign DeltaMath
instead, it replaces this page, you strike through the score cell on the cover, and this page
stays in the packet as extra practice. Never assign both.
\end{skillbox}

\boxguard
\begin{skillbox}[Watch For (while circulating)]{redbox}
\begin{itemize}[leftmargin=1.2em, nosep]
  \item \textbf{``$r$ is close to $\pm1$, so the line is appropriate''} (notes~7, HW~B1,
        AP~MC1 --- the target misconception). Probe: \emph{``Read me the residual plot, not
        $r$. Is there a pattern?''}
  \item \textbf{The reverse error: ``$r$ is only 0.80, so a line is wrong''} (notes~11).
        Probe: \emph{``Show me the curve.''}
  \item \textbf{``Far from the line means influential''} (notes~3, HW~A2). Probe:
        \emph{``Where is it left to right? Now read the slope with and without it.''}
  \item \textbf{Calling every unusual point an outlier} (notes~1--2, 12; AP~(a)). Probe:
        \emph{``Far from the trend, or far out in $x$? Those are different words.''}
  \item \textbf{Stopping at $\widehat{\log y}$} (notes~9, HW~B2, AP~MC2): reporting 0.71 mg/L
        or \$2.24. Probe: \emph{``Is 0.71 a blood level or a log? Undo it.''}
\end{itemize}
Cold-call prompts: \emph{``Give me a strong association that a line should not model.''} \quad
\emph{``Where would you put one point to change the slope the most?''}
\end{skillbox}

\boxguard
\begin{skillbox}[Close \& Assign (8 min)]{goldbox}
\textbf{Students start the homework here, in class, alone and silent --- this is the individual
practice, and the first minutes of it are your best formative read.} Hand the launch first ---
five items on paper, scored, due the first class after two study halls --- and point out that
the AP Practice page is theirs to work before the unit test. Circulate on \textbf{Part~B
item~1}. Sort what you see into three piles as you walk --- disagreed and named the U (ready);
disagreed with no reason, or ``because $r$ isn't 1'' (ready, needs the language); agreed with
the classmate (the misconception survived) --- and open the unit review by reteaching the
residual-plot row if that third pile ran past a third of the room. Then close on what changed:
\emph{``You walked in trusting $r$. You walk out checking the residual plot first.''}
\textbf{Preview:} this is the last lesson of Unit~2 --- next comes the unit test (sample test
at the back of the unit packet), then Unit~3, \emph{Collecting Data}, where association meets
the question of cause.
\end{skillbox}

% ── Teacher notes — one per component, in packet order (four: there is no activity) ──
\begin{teachernote}[Warm-Up]
Item~3 is the seed and item~2 the quiet setup. Give four minutes total. On item~3 many students
will choose P because it looks further from the line; do not settle it --- write both answers on
the board and say problem~3 decides it with numbers. On item~2, make someone say \emph{why} B
rules out a line; that sentence comes back word for word at problem~7, when an impressive $r$
tries to overrule it. Item~1 is a 60-second spiral.
\end{teachernote}

\begin{teachernote}[Guided Notes \& Practice]
\textbf{Pace it 24 / 10, then 8 for the debrief, then 8 for the homework start.} Three rows, 13
problems. \textbf{You work 1, 5, and 8}; the class works the rest. Unusual Points must move ---
ten minutes --- because the time belongs to the Residual Plots row, which carries the crux.
\textbf{Problem~3 is the trap} (A is far off the trend and barely moves the slope);
\textbf{problem~7 is the crux} ($r = -0.94$ beside a U). Take a hand count on 7 before any
explanation, and do not resolve it faster than the room can follow: the callout sentence is the
resolution. The transformation is deliberately light --- read from output, one prediction; if
time is short, do problem~9 yourself and keep problem~8. \textbf{The Lemonade Stand is where
the pen changes hands}; problem~11 is the crux in reverse and the one to protect. If you only
get through 10, 11, and 12, that is the right trade. \textbf{Your second formative read is
the homework start}: Part~B item~1, three piles, and reteach the Residual Plots row in the
unit review if the third pile passes a third of the room.
\textbf{Debrief (8 min), spoken}, opening on the almost-right problem~7. \textbf{Do not let the
debrief eat the last eight minutes} --- the homework start is not optional padding.
\end{teachernote}

\begin{teachernote}[AP Practice]
\textbf{This page is not required and not scored} --- the graded practice is the homework
behind it. It doubles as unit-test review, so go over it in the review class. \textbf{MC~1 is
the one worth reading closely}: (A) is the target misconception verbatim; if more than a handful
choose it, return to the antibiotic residual plot before the test. On the free response, part
(b) must quote numbers from the output --- ``the slope changes'' without $-0.103$ and $-0.055$
is incomplete --- and part (c) must say \emph{predicted} price and carry units.
\end{teachernote}

\begin{teachernote}[Homework]
\textbf{Scored, 10 points (2 per item), due the first class after two study halls.} The eight
minutes at the end of the period are a \emph{start}: put students on B1 first while you can
still catch a wrong start, then A1. \textbf{B1 is the item to read first when scoring} --- a
student who agrees with the classmate has kept the target misconception. On A2, (C) without a
justification earns 1 of 2; the justification must say the point keeps the slope because it is
\emph{on} the line. On B2, a final answer of 2.24 (never undoing the log) earns 1 of 2.
\textbf{Paper is the default.} On the occasions you assign DeltaMath in its place, strike the
score cell on the cover and tell students this page is now optional practice. Do not assign
both.
\end{teachernote}
\end{document}
